\section{Overview}\label{sec:overview}

\textbf{Spectre} is a spectrogram visualizer for VCV Rack 2. It shows how a
signal's frequency content changes over time: frequency runs from bottom to
top, successive spectra occupy columns across the display, and color
represents magnitude. Spectre has one signal input and no audio outputs.
Patch it alongside your audio path to inspect a signal without changing what
you hear.

Fourier and Spectre belong to the Arhythmetic Units Fourier plugin. Use
Fourier to compare the current spectra of up to four signals, and Spectre to
follow one signal's spectrum over time. Both provide the same window choices
and smoothing controls, but their display axes, input-gain defaults, and Run
behavior differ as described below.

Spectre uses a short-time Fourier transform (STFT) with a fixed FFT length of
2048 samples, a hop size of 1024 samples (50\% overlap), and a history of 512
spectra. FFT length, hop size, and history length are not adjustable. You can
choose the window function, time and frequency smoothing, visible frequency
range, display slope, intensity range, and color map.

At sample rate $f_s$, frequency bins are spaced $f_s/2048$ Hz apart, each
analysis window spans $2048/f_s$ seconds, and a complete display sweep takes
$512 \times 1024/f_s$ seconds. At $48\,\mathrm{kHz}$, these are approximately
$23.4\,\mathrm{Hz}$, $42.7\,\mathrm{ms}$, and $10.92\,\mathrm{s}$
respectively. At $44.1\,\mathrm{kHz}$, a sweep takes about
$11.89\,\mathrm{s}$. Increasing Rack's sample rate shortens the history
duration and increases the spacing between bins.
